\section{\filetotitle}
Algorithms for simulating nonadiabatic dynamics are often compared to the exact solution of the \acrshort{tdse} for model systems. This is a common practice in the field of nonadiabatic dynamics, as it allows researchers to benchmark their methods against a known solution. Having a program that can solve the \acrshort{tdse} for a given system is essential for this purpose. In this chapter, we will discuss some of the challenges and benchmarking of quantum dynamics simulations. We will also discuss some tricks and techniques that can be used to improve the efficiency and accuracy of these simulations. Finally, we will discuss some of the applications of quantum dynamics simulations and how they can be used to gain insights into the behavior of quantum systems.

The method used for all simulations in this work is the Split-Operator method, described in detail in Sec.~\ref{sec:split_operator_method}. The algorithm was, as part of this work, implemented in a custom-built quantum chemistry software package called \textsc{Zinq}.\supercite{10.5281/zenodo.18386143} The implementation is designed to allow for simulations of arbitrary dimensions and states, making it a versatile tool for studying a wide range of quantum systems. The software is also optimized for memory efficiency instead of speed, which allows for simulations of larger systems than would otherwise be impossible. In the following sections, we will look at the computational framework and algorithmic design of the implementation in more detail, and discuss some of the challenges and benchmarking of quantum dynamics simulations.

\subsection{Algorithmic Design for Arbitrary Dimensions and States}
While writing a program to propagate a wavefunction in harmonic oscillator using the Split-Operator method is relatively straightforward, extending the implementation to allow for arbitrary dimensions and states is more challenging. One of the main challenges is that the number of grid points required to represent the wavefunction increases exponentially with the number of dimensions. This can lead to memory issues, as well as increased computational time.

For illustration, let's analyze the memory requirements for $K$-dimensional $L$-state simulations with $N$ grid points per dimension using a trivial algorithm. We, of course, need to store the wavefunction, which requires $LN^K$ complex numbers. A sensible thing to do is to store the kinetic and potential energy operators as well, since they do not change for time-independent Hamiltonians. Since these propagators are $L\times L$ matrices in diabatic basis and we need to store them for each grid point, this requires $2L^2N^K$ complex numbers. Finally, we do not want to transform the wavefunction to and from the momentum space every time we want to calculate the kinetic energy, so we also need to store the Fourier transform of the wavefunction, which requires another $LN^K$ complex numbers. This leaves us with a total memory requirement of $2LN^K(L+1)$ complex numbers. One might think that the computer can handle this, but for a 3D 2-state simulation with $N=512$ grid points per dimension (which is pretty standard number of points for lower dimensions) we already need over 25~GB of memory (assuming 16 bytes per complex number).

This, of course, is the na\"ive approach, and there are several tricks that can be used to reduce the memory requirements. For example, instead of storing the kinetic and potential energy operators, we can calculate them on the fly. This reduces the memory requirements to $2LN^K$ complex numbers, which is a significant improvement. However, this comes at the cost of increased computational time, as we need to calculate the operators at each time step, but the on-the-fly calculation of the operators is still necessary for time-dependent Hamiltonians anyway. Another trick is to perform the Fourier transform of the wavefunction in-place, which reduces the memory requirements to $LN^K$ complex numbers. This, however, is not implemented in \textsc{Zinq} as it could introduce numerical errors due to the need of inverse Fourier transform to get back to the position space after e.g. calculating the kinetic energy in the momentum space. Instead, we perform the Fourier transform out-of-place, but state-wise, which allows us to reuse the same memory for transforming several states of the wavefunction at a time, which is a good compromise between memory efficiency and numerical accuracy. With these tricks, we reduced the initial memory requirement of $2LN^K(L+1)$ complex numbers to $(L+1)N^K$ complex numbers, which is a significant improvement and allows us to perform simulations of larger systems and, among other things, allowed us to perform a large number of 2D simulations required for the work described in Sec.~\ref{sec:spin_orbit_induced_charge_transfer}. For comparison, the 3D 2-state simulation with $N=512$ grid points per dimension now only requires about 6~GB of memory, which is a significant improvement over the initial 25~GB.

Another remark that a potential quantum dynamics practitioner might have is how to implement arbitrary dimensions in a compiled language like C, C++, Zig or Fortran. The problem is that the number of dimensions is not known at compile time, which makes it highly impractical and memory inefficient to use multi-dimensional arrays, since they require the number of dimensions to be known at compile time. One solution to this problem is to use a single-dimensional array to store the wavefunction, and then use a mapping function to convert between the multi-dimensional indices and the single-dimensional index. This allows us to implement arbitrary dimensions without the need for multi-dimensional arrays, and is the approach used in \textsc{Zinq}.

Last thing to mention is the extension to arbitrary number of states. Besides some minor problems with memory management, we need to calculate the exponential of the potential energy matrix, as seen in Eq.~\ref{eq:split_operator_propagation}. While the exponential of a 2x2 matrix can be calculated analytically, we need to use numerical methods to calculate the exponential of larger matrices. In \textsc{Zinq}, we use the spectral decomposition of the potential energy matrix to calculate its exponential. This is as simple as finding diagonal matrix $\Lambda$ and unitary matrix $U$ such that $\tilde{\symbf{V}}=\symbf{U}\symbf{\Lambda}\symbf{U}^\dagger$, which allows us to calculate the exponential as $\e^{-\i\tilde{\symbf{V}}\Delta t}=\symbf{U}\e^{-\i\symbf{\Lambda}\Delta t}\symbf{U}^\dagger$. Where the exponential of the diagonal matrix $\Lambda$ is simply the exponential of its diagonal elements.
\subsection{Gauge Invariance of Adiabatic Eigenvectors}
As mentioned in the previous section, some operations in the Split-Operator method, such as the calculation of the exponential of the potential energy matrix or extracting adiabatic wavefunction from the diabatic wavefunction, require the diagonalization of the diabatic potential energy matrix $\tilde{\symbf{V}}(\symbf{R})$ at each spatial grid point $\symbf{R}_i$. This diagonalization yields a set of adiabatic eigenvectors $\symbf{U}(\symbf{R}_i)$ and the potential in the adiabatic representation $\symbf{\Lambda}(\symbf{R}_i)$.

The issue with this transformation is that the adiabatic eigenvectors are not unique, and can be multiplied by an arbitrary phase factor $\e^{\i\phi(\symbf{R}_i)}$ without changing the physics of the system. This is known as gauge invariance. For strictly real symmetric potential energy matrices, the gauge freedom is reduced to a sign change. While this gauge invariance does not affect the physical observables, it can cause some problems while analyzing the results of the simulations. For example, should we want to analyze our wavefunction in the adiabatic representation, we need to be careful when comparing the adiabatic wavefunction at different grid points, as the gauge choice can lead to discontinuities in the wavefunction. Another problem arises when we want to extract coherence between different states from the wavefunction.

When we project our diabatic wavefunction onto the adiabatic basis at a specific grid point $\symbf{R}_i$, we get a set of adiabatic wavefunction coefficients, now only for two-state system, $c_1(\symbf{R}_i)$ and $c_2(\symbf{R}_i)$. The electronic coherence between the two states at this grid point is then
\begin{equation}
\rho_{12}(\symbf{R}_i)=c_1(\symbf{R}_i)c_2^\ast(\symbf{R}_i).
\end{equation}
Now, since the adiabatic eigenvectors are not unique, we can arbitrarily multiply the coefficients by a phase factor $\e^{\i\phi(\symbf{R}_i)}$ (or $-1$ for real symmetric potential energy matrices), which would create a discontinuity in the coherence $\rho_{12}(\symbf{R}_i)$, even though the physical observables remain unchanged. This does not only make the analysis of the results more difficult, but can also lead to incorrect macroscopic coherence values if we want to integrate the coherence over the spatial grid. In this current work, this problem had to be solved due to the coherence analysis performed in the work described in Sec.~\ref{sec:coherence_in_mixed_qc_dynamics}.

The usual way to solve this problem is to choose a gauge that is smooth across the spatial grid. This can be done by tracking the gauge sequentially. Starting from a boundary, the eigenvectors at grid point $\symbf{R}_i$ are compared to the already fixed eigenvectors at the adjacent grid point $\symbf{R}_{i−1}$. We calculate the overlap between the adjacent eigenvectors, and use this overlap to determine the phase factor that we need to apply to the adiabatic eigenvectors at this grid point to make them smooth with respect to the reference gauge. For real symmetric potential energy matrices, this is simply a matter of checking the sign of the overlap and multiplying the adiabatic eigenvectors by $-1$ if the overlap is negative. This way, we can ensure that the adiabatic eigenvectors are smooth across the spatial grid, which allows us to analyze the results of the simulations without worrying about gauge discontinuities.
\subsection{Boundary Reflection and Complex Absorbing Potentials}
A second major problem that arises when discretizing the spatial grid for quantum dynamics simulations is the issue of boundary reflection. Because the wavefunction is propagated using \acrshort{fft}, we implicitly assume periodic boundary conditions, which means that the wavefunction that reaches the boundary of the grid will reenter the grid from the opposite side. This phenomenon also causes other numerical instabilities at the edges, since the potential rarely wraps around. The interaction of the wavefunction with the potential at the edges is therefore undesirable.

One way to solve this problem is, of course, to increase the size of the grid, so that the wavefunction does not reach the boundary during the simulation. However, this is not always possible, especially for higher-dimensional simulations. Another way to solve this problem is to use a \acrfull{cap}. A \acrshort{cap} introduces a strictly imaginary potential which can absorb the wavefunction as it approaches as
\begin{equation}
\symbf{H}(\symbf{R})=\symbf{T}(\symbf{R})+\symbf{V}(\symbf{R})-\i w(\symbf{R})\symbf{I},
\end{equation}
where $w(\symbf{R})$ is the \acrshort{cap}. And since the \acrshort{cap} is usually applied to the diagonal of the potential energy matrix, it is represented as a scalar function and the selection of basis does not matter. The \acrshort{cap} is usually designed to be zero in the region of interest, and then gradually increase as we approach the boundary of the grid. When this non-Hermitian Hamiltonian is used in the Split-Operator method, we can simply split the exponential of the potential energy matrix into the exponential of the real part and the exponential of the imaginary part as
\begin{equation}
\e^{-\i\frac{\Delta t}{2}(\symbf{V}(\symbf{R})-\i w(\symbf{R})\symbf{I})}=\e^{-\i\frac{\Delta t}{2}\symbf{V}(\symbf{R})}\e^{-\frac{\Delta t}{2}w(\symbf{R})\symbf{I}}.
\end{equation}
The real part of the exponent simply acts as a dampening factor, which reduces the amplitude of the wavefunction as it approaches the area where the \acrshort{cap} is non-zero. This creates a finite, memory-efficient grid that mathematically mimics an infinite, open boundary. There are several forms of \acrshort{cap} that can be used, such as the quadratic \acrshort{cap}, the exponential \acrshort{cap}, and the sine-squared \acrshort{cap}. The choice of \acrshort{cap} depends on the specific system being studied and the desired level of accuracy.
\subsection{Benchmarking and Validation}
